\documentclass[12pt,letterpaper, onecolumn]{exam}
\usepackage{amsmath}
\usepackage{amssymb}
\usepackage{amsthm}
\usepackage{graphicx}
\usepackage{mathrsfs}
\usepackage{xcolor}
\usepackage[lmargin=71pt, tmargin=1.2in]{geometry}  %For centering solution box
\usepackage[english]{babel}
\newtheorem{theorem}{Theorem}
\newtheorem{lemma}[theorem]{Lemma}
\usepackage{accents}
\newcommand{\ubar}[1]{\underaccent{\bar}{#1}}
\newcommand{\bs}[1]{\boldsymbol{#1}}
\newcommand{\bm}[1]{\mathbf{#1}}
\newcommand{\der}[2]{\frac{d #1}{d#2}}
\newcommand{\pder}[3]{\frac{\delta^{#3} #1}{\delta^{#3}#2}}
\newcommand{\rs}{$X_1, ..., X_n$}
\newcommand{\nsum}{\sum_{i=1}^n}
\newcommand{\nprod}{\prod_{i=1}^n}
\newcommand*{\Comb}[2]{{}^{#1}C_{#2}}%
\newcommand{\ifc}{\textrm{ if }}
\newcommand{\real}{\mathbb{R}}
\newcommand{\nat}{\mathbb{N}}
\newcommand{\power}{\mathcal{P}}
\newcommand{\rational}{\mathbb{Q}}
\newcommand{\cutelim}{\lim\limits_{n\to\infty}}
\newcommand{\salg}{$\sigma$-algebra }

\newcommand{\calA}{\mathcal{A}}
\newcommand{\calO}{\mathcal{O}}
\newcommand{\calC}{\mathcal{C}}
\newcommand{\calB}{\mathcal{B}}
\newcommand{\calI}{\mathcal{I}}
\newcommand{\calF}{\mathcal{F}}
\newcommand{\calG}{\mathcal{G}}
\newcommand{\gensig}[1]{\sigma\langle #1 \rangle} % notation for a sigma algebra generated by some set 
\newcommand{\om}{\mu^*} % outer measure 
\newcommand{\rat}{\mathbb{Q}}
\newcommand{\tribrac}[1]{\langle #1 \rangle}
\newcommand{\borel}{\mathcal{B}(\mathbb{R})}
\newcommand{\boreltwo}{\mathcal{B}(\mathbb{R}^2)}
\newcommand{\borelk}[1]{\mathcal{B}(\mathbb{R}^#1)}
\newcommand{\probspace}{(\Omega, \calF, P)}

\newcommand{\Iunion}{\bigcup_{\alpha \in I}}
\newcommand{\Iintersect}{\bigcap_{\alpha \in I}}

\newcommand{\dm}{d \mu}
\newcommand{\cutesum}[1]{\sum_{#1 = 1}^\infty}
\newcommand{\lp}{\left (}
\newcommand{\rp}{\right )}

% \chead{\hline} % Un-comment to draw line below header
\thispagestyle{empty}   %For removing header/footer from page 1

\begin{document}

\begingroup  
    \centering
    \LARGE STAT 6410\\
    \LARGE Homework 5\\[0.5em]
    \large Sarah Kilgard\par
\endgroup
\rule{\textwidth}{0.4pt}
\pointsdroppedatright   %Self-explanatory
\printanswers
\renewcommand{\solutiontitle}{\noindent\textbf{Ans:}\enspace}   %Replace "Ans:" with starting keyword in solution box

\begin{questions}
%%%%%%%%%%% Question 1 %%%%%%%%%%%%%%%%%%%%    
    \question Problem 2.15(c) \\
    Consider the probability space $((0, 1), \calB((0, 1)), m)$, where $m(\cdot)$ is the Lebesgue measure. \\
    (c) Let $F:\real \to \real$ be a cdf. For $0 < x < 1$, let 
    \[F_1^{-1}(x) = \inf\{y: y \in \real, F(y) \geq x\}\]
    \[F_2^{-1}(x) = \sup\{y: y \in \real, F(y) \leq x\}\]
    Let $Z_i$ be the random variable defined by $Z_i = F_i^{-1}(x), \: 0 < x < 1, \: i = 1, 2.$
    \begin{itemize}
        \item[(i)] Find the cdf of $Z_i$, $i = 1, 2$ \\
        (\textbf{Hint:} Verify using the right continuity of $F$ that for any $0 < x < 1$, $t \in \real$, $F(t) \geq x \iff F_1^{-1}(x) \leq t.$)
        \item[(ii)] Show also that $F_1^{-1}(\cdot)$ is left continuous and $F_2^{-1}(\cdot)$ is right continuous.  
    \end{itemize}
    \begin{solution}
    (i) 
    
    \textbf{Finding the cdf of $Z_1$} 
    \begin{proof}
        Let $t \in \real$ and $x \in (0, 1)$ be given. 
        
        \textbf{Part 1:} WTS $F(t) \geq x \implies F_1^{-1}(x) \leq t.$ \\
        Suppose $F(t) \geq x$. \\
        $F_1^{-1}(x) = \inf\{y: y \in \real, F(y) \geq x\} $. \\
        From the definition of $F_1^{-1}(\cdot)$, we know that there exists no $r \in \real$ such that $F(r) \geq x$ and $r \leq F_1^{-1}(x)$. \\
        Thus, since $x \leq F(t)$, we know that, $F_1^{-1}(x) \leq t$. 

        \textbf{Part 2:} WTS $F(t) \geq x \impliedby F_1^{-1}(x) \leq t.$ \\
        Suppose $F_1^{-1}(x) \leq t$. \\
        Define $\{y_n\}_{n \geq 1}$ to be a decreasing sequence from $\{y: y \in \real, F(y) \geq x\}$. \\
        By the right continuity of $F(\cdot)$ (F is a cdf), $\cutelim y_n = F_1^{-1}(x)$. 

        $t \geq F_1^{-1}(x) \implies F(t) \geq F(F_1^{-1}(x) ) = F(\cutelim y_n) = \cutelim F(y_n) \geq x$

        $\therefore \: F(t) \geq x \impliedby F_1^{-1}(x) \leq t.$

        $\therefore \: F(t) \geq x \iff F_1^{-1}(x) \leq t.$

        \begin{align*}
            & P(Z_1 \leq t) = P(F_1^{-1}(x) \leq t) \\
            & = P(x \leq F(t)) && \textrm{(since }F(t) \geq x \iff F_1^{-1}(x) \leq t) \\
            & = m(\{x: x \leq F(t)\}) = F(t)
        \end{align*}

        Thus, the cdf of $Z_1$ is $F(t)$ for all $t \in \real$.

        \textbf{Finding the cdf of $Z_2$}

        Using a similar logic, we can show that for all $x \in (0, 1)$ $F(t) > x \: \implies \: F_2^{-1} (x) \leq t$. This means that $F(t) = m(\{x: F(t) > x\}) \leq m(\{F_2^{-1}(x) \leq t\})$. We can, in a very similar manner to the first part of this problem show that $F(t) > x \: \impliedby \: F_2^{-1} (x) \leq t$. This means that $m(\{F_2^{-1}(x) \leq t\}) \leq m(\{x: F(t) > x\}) = F(t)$. Thus, $F(t) = m(\{F_2^{-1}(x) \leq t\})$

        $$P(Z_2 \leq t) = m(\{F_2^{-1}(x) \leq t\}) = F(t)$$

        Thus, the cdf of $Z_2$ is $F(t)$ for all $t \in \real$.
        
        (ii)

        also use the hint in (ii)
        \begin{proof}
            \textbf{Part 1:} WTS $F_1^{-1}(\cdot)$ is left continuous

            Let $\epsilon > 0$ be given. There exists $\delta_\epsilon$ such that $F(F^{-1}(x) - \epsilon) < x - \delta_\epsilon$. From $F(t) \geq x \iff F_1^{-1}(x) \leq t$ (in part (i)), we know that $F^{-1}_1(x - \delta_\epsilon) > F_1^{-1}(x-\epsilon)$. 

            Define $x^* = (x - \delta_\epsilon, x)$. \\
            $F_1^{-1}(x - \delta_\epsilon) <  F_1^{-1}(x*) \implies$ \\
            $F_1^{-1}(x) - F_1^{-1}(x*) < F_1^{-1}(x) - F_1^{-1}(x - \delta_\epsilon) < \epsilon$

            Thus, $F^{-1}_1(\cdot)$ is left continuous. 
            
            \textbf{Part 2:} WTS $F_2^{-1}(\cdot)$ is right continuous

            Again, referring back to (i): for all $\epsilon > 0$, there exists some $\delta_\epsilon > 0$ such that $F(F^{-1}_2(x + \epsilon) > x + \delta_\epsilon \: \implies \: F_2^{-1}(x + \delta_\epsilon) \leq F_2^{-1}(x) + \epsilon$. Thus, $F_2(\cdot)$ is right continuous. 
        \end{proof}
    \end{proof}
    \end{solution}

%%%%%%%%%%% Question 2 %%%%%%%%%%%%%%%%%%%%    
    \question Problem 2.20 \\
    Apply Corollary 2.3.5 to show that for any collection $\{a_{ij}: i, j \in \nat\}$ of nonnegative numbers, 
    \[\sum_{i = 1}^\infty \left (\sum_{j = 1}^\infty a_{ij}\right) = \sum_{j = 1}^\infty \left(\sum_{i = 1}^\infty a_{ij}\right).  \]

    \textbf{Corollary 2.3.5:}
    \textit{Let $\{h_n\}_{n \geq 1}$ be a sequence of nonnegative measurable functions on a measure space $(\Omega, \calF, \mu)$. Then}
    \[\int\left(\sum_{n = 1}^\infty h_n\right) d\mu = \sum_{n = 1}^\infty \int h_n d\mu\]
    \begin{solution}
    \begin{proof}
    Define $\mu(\cdot)$ to be a counting measure. \\
    If $\{b_k: k \in \nat\}$ is a collection of nonnegative numbers:
    \[\cutesum{k}b_k = \cutesum{k} b_k I(\{k\}) = \int b_k d \mu(\{i\})\]
    We can apply this to our problem at hand:
    \begin{align*}
        & \cutesum{i}\left(\cutesum{j} a_{ij}\right) = \cutesum{i}\lp \cutesum{j} a_{ij}\rp I(\{i\}) = \int \lp \cutesum{j} a_{ij}\rp d\mu(\{i\}) \\
        & = \cutesum{j} \lp \int a_{ij} d \mu(\{i\}) \rp && \textrm{(by \textbf{Corollary 2.3.5)}} \\
        & \cutesum{j}\lp \cutesum{i} a_{ij} I(\{i\}) \rp = \cutesum{j} \lp \cutesum{i} a_{ij} \times 1 \rp \\
        & = \cutesum{j} \lp \cutesum{i} a_{ij} \rp
    \end{align*}
    \end{proof}
    \end{solution}

%%%%%%%%%%% Question 3 %%%%%%%%%%%%%%%%%%%%    
    \question Problem 2.25 \\
    If $f(x) = I_{\rat_1}(x)$ where $\rat_1 = \rat \cap [0,1]$, $\rat$ being the set of all rational, then show that for any partition $P,$ $U(P, f) = 1$ and $L(P, f) = 0$. 
    \begin{solution}
        \begin{proof}
            Define a partition, P ,on $\rat_1$ as $P = (x_0 = 0, x_1, ... , x_n = 1)$ \\
            Rationals are dense in the real numbers, thus for all $0 \leq i \leq 1$, there exists a rational number in $[x_i, x_{i+1}]$. \\Thus, $\sup(f(x): x_i \leq x \leq x_{i+1}) = 1$ for all $i.$ \\
            $U(P, f) = \sum_{i = 1}^n \sup(f(x): x_i \leq x \leq x_{i+1}) (x_{i+1} - x_{i}) = \sum_{i = 1}^n 1 (x_{i+1} - x_{i}) = \sum_{i =1}^n (x_{i+1} - x_{i}) = 1 - 0 = 1$.

            Similarly, irrational numbers are dense in the real numbers, thus for all $0 \leq i \leq 1$, there exists an irrational number in $[x_i, x_{i+1}]$. \\Thus, $\inf(f(x): x_i \leq x \leq x_{i+1}) = 0$ for all $i.$ \\
            $L(P, f) = \sum_{i = 1}^n \sup(f(x): x_i \leq x \leq x_{i+1}) (x_{i+1} - x_{i}) = \sum_{i = 1}^n 0 (x_{i+1} - x_{i}) = 0$.
        \end{proof}
    \end{solution}

%%%%%%%%%%% Question 4 %%%%%%%%%%%%%%%%%%%%    
    \question Extra Problem (A) \\
    Let $f: \Omega \to \real$ be a $\tribrac{\calF, \borel}$-measurable \textit{nonnegative} function. Define for $n \geq 1$, and $x \in \Omega$
    \[f_n(x) = \begin{cases}
        \frac{(i - 1)}{2^n}, \ifc \frac{(i - 1)}{2^n} \leq f(x) < \frac{i}{2^n} \textrm{for } i = 1, ... , n2^n \\
        n, \ifc f(x) \geq n
    \end{cases}\]
    Then prove for the following:
    \begin{itemize}
        \item[(i)] For each $n \geq 1$, $f_n(x)$ is $\tribrac{\calF, \: \borel}$-measurable.
        \item[(ii)] For each $x \in \Omega$, $f_n(x)$ is increasing with $n$ (i.e. $\{f_n(\cdot)\}$ is an increasing sequence of functions). 
        \item[(iii)] Prove that $f_n(x) \to f(x)$ pointwise (i.e. for each fixed $x$), as $n \to \infty$ 
    \end{itemize}
    \begin{solution}
        (i) 
        \begin{proof} 
        Let $n \geq 1$ be given. \\
        We can write $f_n(x) = \sum_{i = 1}^{n2^n} \frac{i - 1}{2^n} I(\frac{(i - 1)}{2^n} \leq f(x) < \frac{i}{2^n}) + n I(f(x) \geq n)$ \\
        $ = \sum_{i = 1}^{n2^n} \frac{i - 1}{2^n} I(x : x \in f^{-1}([\frac{i - 1}{2^n}), \frac{i}{2^n}))) + n I(x : x \in f^{-1}([n, \infty)))$ \\
        $[0, \frac{1}{2^n}), [\frac{1}{2^n}, \frac{2}{2^n}), ... , [\frac{n2^n - 1}{2^n}, n), [n, \infty) \in \borel $ \\
        Thus, for all $i \in \{1, 2, ... , n2^n\}$, $\{x : x \in f^{-1}([\frac{i - 1}{2^n}), \frac{i}{2^n}))\} \in \calF$ and \\ $\{x: x \in f^{-1}([n, \infty))\} \in \calF$ \\
        Thus, $f_n(x)$ is a sum of indicator functions multiplied by constants. \\ All of the indicator functions are $\tribrac{\calF, \borel}-$ measurable, since all of the sets by which the indicator function is defined are elements of $\calF$. \\
        By \textbf{Corollary 2.1.4}, $\tribrac{\calF, \borel}-$ measurable functions are closed under addition and scalar multiplication, $f_n(x)$ is $\tribrac{\calF, \borel}-$ measurable.
        \end{proof}

        (ii) 
        \begin{proof}
            Let $x \in \Omega$ be given. \\
            First, let's rewrite $f_n(x)$ to make it easier to work with:
            \begin{align*}
                &f_n(x) = \sum_{i = 1}^{n2^n} \frac{i - 1}{2^n} I(\frac{(i - 1)}{2^n} \leq f(x) < \frac{i}{2^n}) + n I(f(x) \geq n) \\
                & = \sum_{i = 1}^{n2^n} \frac{i - 1}{2^n} I((i - 1) \leq 2^nf(x) < i) + n I(f(x) \geq n) \\
                & = \frac{\lfloor \min(2^n f(x), 2^n n)\rfloor}{2^n}
            \end{align*}

            Now to answering the question at hand: 
            \begin{align*}
                &f_{n+1}(x) = \frac{\lfloor \min(2^{n+1} f(x), 2^{n+1} (n+1))\rfloor}{2^{n+1}} \\
                & = \frac{\lfloor 2\min(2^{n} f(x), 2^{n} (n+1))\rfloor}{2^{n+1}} = \frac{2\lfloor\min(2^{n} f(x), 2^{n} (n+1))\rfloor}{2^{n+1}} \\
                & = \frac{\lfloor\min(2^{n} f(x), 2^{n} (n+1))\rfloor}{2^{n}} \geq \frac{\lfloor\min(2^{n} f(x), 2^{n} (n+1))\rfloor}{2^{n}} \\
                & = f_n(x) 
            \end{align*}
        \end{proof}
        (iii) 
        \begin{proof}
            Let $x \in \Omega$ be given. \\
            Thus, there exists some $N \in \nat$ such that $N \geq f(x)$. Since we are investigating the behavior of n it approaches $\infty$, $\min(2^nf(x), 2^n n) = 2^n f(x)$ as $n \to \infty.$

            \begin{align*}
            & \lim\limits_{n \to \infty} |f_n(x) - f(x)| = \cutelim \left|\frac{2^n f(x) -  \lfloor \min(2^n f(x), 2^n n)\rfloor}{2^n}\right| \\
            & = \cutelim \frac{|2^n f(x) -  \lfloor 2^n f(x)\rfloor|}{2^n} \leq \cutelim \frac{1}{2^n} = 0
            \end{align*}
            Thus, $f_n(x) \to f(x)$ point-wise as $n \to \infty$ for all $x \in \Omega$

           
        \end{proof}
    \end{solution}

%%%%%%%%%%% Question 5 %%%%%%%%%%%%%%%%%%%%    
    \question Extra Problem (B) \\
    If $f, g \in L_1(\Omega, \calF, \mu)$, \\
    show 
    \[\min(f, g) \in L_1(\Omega, \calF, \mu)\]
    and 
    \[\min\left(\int fd\mu, \: \int g d\mu\right) \geq \int \min(f,g)d\mu.\]
    \begin{solution}
        WTS(Want to Show): $\min(f, g) \in L_1(\Omega, \calF, \mu)$
        \begin{proof}
            $|\min(f, g)| \leq |f| + |g|$ \\
            Since, $f, g \in L_1(\Omega, \calF, \mu)$, $\int |f| d \mu < \infty$ and $\int |g| d \mu < \infty$. \\ This means that $\int |f| + |g| d \mu < \infty$. 
            By monotonicity, $\int |\min(f, g)| d \mu \leq \int |f| + |g| d \mu < \infty$. \\ Thus, $\min(f, g) \in L_1(\Omega, \calF, \mu)$
        \end{proof}

        WTS: $\min\left(\int fd\mu, \: \int g d\mu\right) \geq \int \min(f,g)d\mu$
        \begin{proof}
            First, we will state this equality: $\min(x, y) = \frac{1}{2}(x + y - |x - y|)$. 
            \begin{align*}
                & \int \min(f, g) \dm = \frac{1}{2}\int f + g - |f - g| \dm \\
                & = \frac{1}{2}\left(\int f\dm + \int g \dm - \int |f - g| \dm \right) \\
                & \leq \frac{1}{2}\left(\int f\dm + \int g \dm - \left|\int f - g \dm\right|\right) && \textrm{(*)}\\
                & = \frac{1}{2}\left(\int f\dm + \int g \dm - \left|\int f \dm - \int g \dm\right|\right) \\
                & = \min\left(\int f \dm, \int g \dm \right)
            \end{align*}

            (*) Since $f$ and $g$ are measurable, $f - g$ is clearly measurable. Thus $\int|f - g| \dm \geq \left|\int f - g \dm \right |$
        \end{proof}
    \end{solution}

\end{questions}
\end{document}